\section{Convolutional Neural Networks (CNNs)}\label{sec:convolutional-neural-networks}

In the previous section, we discussed \textbf{image classification} using relatively simple models:
\begin{itemize}
    \item Linear classifiers, which are based on linear decision boundaries
    \item Multi-layer perceptrons, which are based on fully connected layers
\end{itemize}
These models can classify images, but they have an important limitation: \textbf{they treat an image simply as a long vector of numbers}. For example, a $32 \times 32$ RGB image becomes:
\begin{equation*}
    32 \times 32 \times 3 = 3072
\end{equation*}
Independent input features. The model does \textbf{not know} that nearby pixels are related, or that edges form shapes, or that shapes form objects. Every pixel is treated almost independently.

\highspace
CNNs were invented to exploit one fundamental property of images: \textbf{pixels that are close together \emph{usually} belong to the same visual structure}. Instead of learning from isolated pixels, CNNs learn \textbf{local visual patterns}, such as: edges, corners, textures, object parts, and complete objects. These patterns are called \textbf{features}, and learning them automatically is the central idea of this section.